\begin{textsample}{POS Dim 2 – ap – Score 28.00 – ap/ap\_em\_12.txt}  \label{ex:f2_pos_084}
EDUCAÇÃO ESCOLAR QUILOMBOLA A luta quilombola por uma educação de qualidade, que preserve sua cultura e que respeite seus aspectos históricos é de longa data, desde a época das reinvindicações pela conquista de terras.
A época do Império, a população negra era impedida de \textbf{frequentar} a escola formal, a qual era restrita por Lei, aos cidadãos brasileiros ( \textbf{art}. 6, item 1 da Constituição de 1824 ).
O decreto 1.331a de 17 de \textbf{fevereiro} de 1854 foi \textbf{instituído} a obrigatoriedade da escola primária para crianças maiores de 07 anos e a gratuidade das escolas primárias e secundárias da corte.
Porém, vale ressaltar que esse decreto explicitamente comprova a ideologia da interdição, ou seja, nas escolas públicas não seriam admitidas crianças com moléstias contagiosas e nem escravas[^1 ] ; assim como também, não havia previsão de instrução para adultos.
A abolição da escravatura no Brasil não livrou os negros africanos e os afro-brasileiros da discriminação racial e suas consequências nefastas.
Mais do que isso, ela passou a ser um dos determinantes do destino social, econômico, político e cultural dos afro-brasileiros ( HASENBALG, 1979 ; SANTOS, 1997 ).
Dessa forma, os negros libertos perceberam que a luta pela liberdade fora apenas o primeiro passo para a obtenção da igualdade no pós-abolição.
Contudo, passaram a lutar pela “ segunda abolição ” ( BASTIDE e FERNANDES, 1955 ) o que poderia proporcionar \textbf{mobilidade} social vertical, ou seja, a busca pela formação educacional para ascender socialmente.
Historicamente, é possível constatar a ausência de leis contundentes em prol da Comunidade Negra no Brasil.
Numa perspectiva mais abrangente da \textbf{garantia} de direitos do cidadão, a Constituição Federal de 1988 inova ao estabelecer que : > Art. 215.
O Estado garantirá a todos o pleno exercício dos direitos culturais e acesso às fontes da cultura nacional, e apoiará e incentivará a valorização e a difusão das manifestações culturais. > § 1º O Estado protegerá as manifestações das culturas populares, indígenas e afro-brasileiras, e das de outros grupos participantes do processo civilizatório nacional. > > Art. 216.
Constitui patrimônio cultural brasileiro os bens de natureza material e imaterial, tomados individualmente ou em conjunto, portadores de referência à identidade, à ação, à memória dos diferentes grupos formadores da sociedade brasileira, nos quais se incluem : > I - as formas de expressão ; > II - os modos de criar, fazer e viver ; > III - as criações científicas, artísticas e tecnológicas ; > IV - as obras, objetos, documentos, edificações e demais espaços \textbf{destinados} às manifestações artístico-culturais ; > § 5º Ficam tombados todos os documentos e os sítios detentores de reminiscências históricas dos antigos quilombos ( BRASIL, 1988 ).
Ainda, na Constituição Federal de 1988, no Ato das Disposições Constitucionais Transitórias, o \textbf{artigo} 68 especifica que “ aos remanescentes das comunidades dos quilombos que estejam ocupando suas terras é reconhecida a propriedade definitiva, devendo o Estado emitir -lhes os títulos respectivos ” ( BRASIL, 1988 ).
O decreto nº 4887, de 20 de \textbf{novembro} de 2003, da Presidência da República, estabelece que : > Art. 1º Os procedimentos administrativos para a identificação, o reconhecimento, a delimitação, a demarcação e a titulação da propriedade definitiva das terras ocupadas por remanescentes das comunidades dos quilombos, de que trata o \textbf{art}. 68 do Ato das Disposições Constitucionais Transitórias, serão procedidos de acordo com o estabelecido neste Decreto. ( BRASIL, 2003 ).
Para finalizar a ação juntamente com o decreto citado, fora \textbf{instituída} a Política Nacional de Promoção da Igualdade Racial ( Decreto nº 4886/2003 ), que apresenta como objetivo principal “ reduzir as desigualdades raciais no Brasil, com ênfase na população negra ”. ( BRASIL, 2003 ).
Contemplar as comunidades quilombolas nas leis, documentos e programas governamentais torna -se imperativo para a \textbf{garantia} de direitos desses povos que foram historicamente excluídos.
Em termos de \textbf{políticas} educacionais, podemos destacar o \textbf{artigo} 26‐A da Lei de \textbf{Diretrizes} e Bases da Educação Nacional ( LDBEN n. 9.394/96 ), introduzido pela Lei nº 10.639/2003, que trata da obrigatoriedade do estudo da História da África e da Cultura afro-brasileira e africana e do ensino das relações étnico-raciais, \textbf{instituindo} o estudo das comunidades remanescentes de quilombos e das experiências negras constituintes da cultura brasileira : > § 1º O conteúdo programático a que se refere o caput deste \textbf{artigo} incluirá o estudo da História da África e dos Africanos, a luta dos negros no Brasil, a cultura negra brasileira e o negro na formação da sociedade nacional, resgatando a contribuição do povo negro nas áreas social, econômica e \textbf{políticas} pertinentes à História do Brasil. > § 2º Os conteúdos referentes à História e Cultura Afro-Brasileira serão ministrados no âmbito de todo o \textbf{currículo} escolar, em especial nas áreas de Educação Artística e de Literatura e História Brasileira. ( BRASIL, 2003 ).
Pelo \textbf{Parecer} CNE/CP nº 03/2004 todo sistema de ensino precisará providenciar “ Registro da história não contada dos negros brasileiros, tais como os remanescentes de quilombos, comunidades e territórios negros urbanos e rurais ” ( BRASIL, 2003, p.9 ).
A Conferência Nacional de Educação ( CONAE ), ocorrida em 2010, incluiu a discussão sobre o direito a educação das comunidades quilombolas.
Como fruto dessas discussões, o \textbf{Parecer} CNE/CEB 07/2010 e a Resolução CNE/CEB 04/2010 que \textbf{instituem} as \textbf{Diretrizes} Curriculares Gerais para a Educação Básica, incluem a educação escolar quilombola como modalidade da educação básica.
Desta forma, somente em 2012 foram dispostas as \textbf{Diretrizes} Curriculares Nacionais para a Educação Escolar Quilombola na Educação Básica por meio da Resolução Nº 08/ 2012 do Conselho Nacional de Educação.
Entende -se que as \textbf{Diretrizes} Curriculares Nacionais para a Educação Escolar Quilombola, representa uma vitória dos movimentos sociais, pois elas nasceram na base, a partir da luta da população negra, mais especificamente do movimento quilombola.
Uma revolução no ensino brasileiro tendo em vista que as referidas \textbf{diretrizes} orientam os sistemas de ensino a valoriza os saberes, as tradições e o patrimônio cultural das comunidades remanescente de quilombos, algo impensável em outras épocas.
A proposta da educação nacional acentua o reconhecimento à diversidade de identidades na valorização do ser humano, e considera no processo, o direito de ser aceito nas especificidades que compõem a nação brasileira.
Nesse sentido, percebe -se que a diferença precisa ser levada em consideração em todos os contextos, sendo a escola um dos espaços importantes para colaborar na superação de todas as formas de discriminação e racismo.
Sendo assim, as propostas pedagógicas, devem acolher com autonomia e o princípio da identidade pessoal e coletiva dos professores, dos estudantes e de todos que convivem nesse espaço social.
Nesse sentido foi criada a Resolução 08/2012 – CNE definindo que : > Art. 9º A Educação Escolar Quilombola compreende : > I - Escolas quilombolas ; > II- Escolas que atendem estudantes oriundos de território quilombola. > Parágrafo Único Entende -se por escola quilombola aquela localizada em território quilombola.
A educação escolar quilombola deve ter como referência valores sociais, culturais, históricos e econômicos dessas comunidades.
Para tal, a escola deverá se tornar um espaço educativo que efetive o diálogo entre o conhecimento escolar e a realidade local, valorize o desenvolvimento sustentável, o trabalho, a cultura, a luta pelo direito à terra e ao território.
Portanto, a escola precisa de \textbf{currículo}, projeto político-pedagógico, espaços, tempos, \textbf{calendários} e temas adequados às características de cada comunidade quilombola para que o direito à diversidade se concretize.
Essa discussão precisa fazer parte da formação inicial e continuada dos professores.
Essa necessidade é enfatizada no Título II dos princípios da educação escolar quilombola : > Art. 7º A Educação Escolar Quilombola rege -se nas suas práticas e ações político-pedagógicas pelos seguintes princípios : > I - direito à igualdade, liberdade, diversidade e pluralidade ; Assim como também no \textbf{artigo} oitavo : > Art. 8º Os princípios da Educação Escolar Quilombola deverão ser garantidos por meio das seguintes ações : > I - construção de escolas públicas em territórios quilombolas, por parte do poder público, sem prejuízo da ação de ONG e outras \textbf{instituições} comunitárias ; > II - adequação da estrutura física das escolas ao contexto quilombola, considerando os aspectos ambientais, econômicos e socioeducacionais de cada quilombo ; > III - \textbf{garantia} de condições de acessibilidade nas escolas ; > IV - presença preferencial de professores e gestores quilombolas nas escolas quilombolas e nas escolas que recebem estudantes oriundos de territórios quilombolas ; > V - \textbf{garantia} de formação inicial e continuada para os docentes para atuação na Educação Escolar Quilombola ; > VI - \textbf{garantia} do protagonismo dos estudantes quilombolas nos processos político-pedagógicos em todas as etapas e modalidades ; > VII - implementação de um \textbf{currículo} escolar aberto, \textbf{flexível} e de caráter interdisciplinar, elaborado de modo a articular o conhecimento escolar e os conhecimentos construídos pelas comunidades quilombolas ; > VIII - implementação de um projeto político-pedagógico que considere as especificidades históricas, culturais, sociais, políticas, econômicas e identitárias das comunidades quilombolas ; > IX - efetivação da \textbf{gestão} democrática da escola com a participação das comunidades quilombolas e suas lideranças ; > X - \textbf{garantia} de alimentação escolar voltada para as especificidades socioculturais das comunidades quilombolas ; > XI - inserção da realidade quilombola em todo o material didático e de apoio pedagógico produzido em articulação com a comunidade, sistemas de ensino e \textbf{instituições} de Educação Superior ; > XII - \textbf{garantia} do ensino de História e Cultura Afro-Brasileira, Africana e Indígena, nos termos da Lei nº 9394/96, com a redação dada pelas Leis nº 10.639/2003 e nº 11.645/2008, e na Resolução CNE/CP nº 1/2004, fundamentada no \textbf{Parecer} CNE/CP nº 3/2004 ; > XIII - efetivação de uma educação escolar voltada para o etnodesenvolvimento e para o desenvolvimento sustentável das comunidades quilombolas ; > XIV - realização de processo educativo escolar que respeite as tradições e o patrimônio cultural dos povos quilombolas ; > XV - \textbf{garantia} da participação dos quilombolas por meio de suas representações próprias em todos os órgãos e espaços deliberativos, consultivos e de monitoramento da \textbf{política} pública e demais temas de seu interesse imediato, conforme reza a Convenção 169 da OIT ; > XVI - articulação da Educação Escolar Quilombola com as demais \textbf{políticas} públicas relacionadas aos direitos dos povos e comunidades tradicionais nas diferentes esferas de \textbf{governo}.
Como Modalidade de Ensino da Educação Básica, a Educação Escolar Quilombola deve ser \textbf{ofertada} em todos os seus seguimentos, de acordo com os assim como, enfatiza a importância da construção dos respectivos Projetos Políticos Pedagógicos Quilombolas.
Resolução CEE/AP nº 075/2009 – Estabelece normas complementares das \textbf{Diretrizes} Curriculares Nacionais.
Inclui a obrigatoriedade da História e Cultura Afro-brasileira e Indígena no \textbf{Currículo} da Educação Básica e Superior no Sistema Estadual de Ensino do Estado do Amapá.
A Resolução nº. 051/2012- CEE/AP estabelece normas complementares às \textbf{diretrizes} curriculares nacionais para a educação das relações étnico-raciais e para o ensino de história e cultura afro-brasileira, africana e indígena no \textbf{currículo} da educação básica e superior no sistema estadual de ensino do estado do Amapá.
Dentre as suas orientações vale ressaltar que : > Art. 4º.
Na vivência da Interdisciplinaridade, as escolas deverão ter presente que a prática da transversalidade valoriza e orienta as atitudes dos estudantes, para a reflexão e análise dos elementos da cultura e dos acontecimentos que ocorrem no contexto social, e : > I – os conteúdos referentes à História e Cultura Afro-Brasileira e Indígena serão ministrados no âmbito de todo o \textbf{currículo} escolar, em especial, nas Disciplinas de Artes, Literatura, História, Geografia e Língua Portuguesa ; > II – o ensino deve ir além da descrição dos acontecimentos e deve procurar desenvolver nos estudantes a capacidade de reconhecer e valorizar a história, a cultura, a identidade e as contribuições dos Afrodescendentes e dos Indígenas na construção histórica, social, econômica e no desenvolvimento da Nação Brasileira ; > III – os conteúdos programáticos devem estar fundamentados em dimensões históricas, sociais, políticas, econômicas, estéticas, religiosas, culturais e antropológicas, referentes à realidade brasileira, em especial a amapaense, destacando as particularidades da história regional, com vistas a combater o racismo e as discriminações que atingem particularmente os afrodescendentes e os indígenas ; > IV – a abordagem temática deve visar a formação de atitudes, posturas e valores que eduquem cidadãos orgulhosos de seu pertencimento étnico-racial, como descendentes de africanos e indígenas ; > V – os conteúdos multidisciplinares devem trabalhar a cultura negra e indígena brasileira, dando destaque aos acontecimentos e realizações próprios da região Norte, em especial do Estado do Amapá. > VI – o centro das abordagens temáticas subsidiadas por recursos didáticos diversos, inclusive pela Pedagogia de Projetos, visa : > a ) a pesquisa, a produção, a leitura, os estudos e a reflexão sobre a temática indígena e africana ; > b ) a implementação de Políticas de Ações Afirmativas que impliquem em justiça e igualdade de direitos sociais, civis, culturais e econômicos. > Art. 9º.
Cada \textbf{instituição} de ensino, no âmbito do Sistema Estadual registrará no requerimento da matrícula de cada aluno, seu pertencimento étnico-racial, garantindo -se o registro da autodeclaração. > Art. 10.
Os estabelecimentos de ensino desenvolverão suas propostas pedagógicas para Educação das Relações Étnico-raciais e o Ensino de História e Cultura Afro-brasileira, Africana e Indígena, elaboradas no âmbito da sua autonomia, no projeto político pedagógico dessas \textbf{Instituições}, obedecendo às recomendações legais. > Art. 11.
O Regimento Escolar deve contemplar normas para avaliação e encaminhamentos de soluções para situações de discriminação, \textbf{prevendo} a adoção de ações didáticas educativas voltadas para o reconhecimento, valorização e respeito à diversidade. > Art. 12.
A proposta pedagógica e o \textbf{regimento} escolar das \textbf{Instituições} de Ensino deverão incluir a educação das relações étnico-raciais, envolvendo toda a comunidade escolar no desenvolvimento dos valores humanos, do respeito aos diferentes biótipos, às manifestações culturais, hábitos e costumes. > Art. 13.
As \textbf{Instituições} de Ensino poderão estabelecer parcerias com grupos culturais do Movimento Negro e Indígena, \textbf{Instituições} Formadoras de Professores, Núcleos de Estudos e Pesquisas, Antropólogos e Sociólogos, com a finalidade de buscar subsídios para planos institucionais, propostas pedagógicas e projetos de ensino. > Art. 14.
Cada \textbf{instituição} de ensino deverá compor \textbf{equipe} interdisciplinar que estará responsável pela supervisão e desenvolvimento de ações que deem conta da aplicação efetiva das \textbf{diretrizes} estabelecidas por esta Resolução ao longo do período letivo e não apenas em datas festivas, pontuais, deslocadas do cotidiano da escola. > Parágrafo único- As \textbf{Instituições} de Ensino, tanto públicas como privadas, deverão contribuir para a preservação da memória das ações desenvolvidas no cumprimento do que preceitua a presente Resolução, por meio de registros diversos.
O Estado do Amapá possui 26 ( vinte e seis ) escolas quilombolas : As comunidades remanescentes de quilombos possuem dimensões educacionais, sociais, políticas e culturais significativas, com particularidades no contexto geográfico e histórico brasileiro, tanto no que diz respeito à localização, quanto à origem.
Dessa forma, as escolas quilombolas desenvolverão suas atividades de acordo com seu projeto político pedagógico.
Construído coletivamente com a comunidade escolar, valorizando os conhecimentos tradicionais, a oralidade, a ancestralidade e a história de cada Comunidade.
Portanto o Núcleo de Educação Étnico Racial – NEER/CEESP/SEED é o órgão responsável pelo assessoramento técnico/pedagógico dos profissionais da educação que desenvolvem os seus trabalhos educacionais nas escolas quilombolas do Estado do Amapá.
Realizando anualmente o Encontro de Gestores Quilombolas que permite a socialização dos problemas encontrados pelas escolas e o aprimoramento do conhecimento abordando temas relacionados a Educação Escolar Quilombola em sua estruturação sistêmica e organizacional.
Observa -se o Estado brasileiro mudando sua postura mediante aos quilombolas, pois passam de uma visão apenas cultural para uma responsabilidade pública na questão política educacional.
Essas conquistas são uma consequência positiva da luta pelo reconhecimento dos direitos coletivos que os quilombolas buscam.
Essas \textbf{diretrizes} levam em consideração a produção cultural, política, econômica e social, fortalecendo a identidade étnico/cultural dos quilombolas do Estado do Amapá através da educação didático/pedagógico desenvolvidas nas unidades de ensino localizadas nessas Comunidades. [ ^1 ] : ROMÃO, Jeruse. ( Org ). “ História da Educação do Negro e Outras História ”.
SECAD – Brasília : Ministério da Educação. 2005. ( Coleção Educação para Todos )

% matched lemmas: art, artigo, calendário, currículo, destinar, diretriz, equipe, fevereiro, flexível, frequentar, garantia, gestão, governo, instituir, instituição, mobilidade, novembro, ofertar, parecer, política, prever, regimento
\end{textsample}
